% Table: Case study — teacher vs student reasoning traces.
\begin{table}[t]
\centering
\caption{Representative case study comparing teacher (Qwen3-32B + GRPO main,
  checkpoint used for private eval) and student (Qwen3-8B + RKL distill) on the
  \emph{same} middle-school prompt (aligned by hash of the user message in
  \texttt{output/eval/*\_middle.jsonl}).  Both assign full credit; the student
  uses finer step-block granularity (more blocks/tokens).  Full traces are in
  Appendix~\ref{app:case_study}.}
\label{tab:case_study}
\vspace{2pt}
\resizebox{\textwidth}{!}{%
\begin{tabular}{@{}p{2.2cm} p{5.5cm} p{5.5cm}@{}}
\toprule
& \textbf{Teacher (32B)} & \textbf{Student (8B)} \\
\midrule
Problem & \multicolumn{2}{p{11cm}}{%
  Quadratic equation $x^2-6x-k^2=0$: (1) prove two distinct real roots;
  (2) given roots $x_1,x_2$ with $x_1+2x_2=14$, find roots and $k$.} \\
\midrule
\# Steps & 2 & 4 \\
\# Tokens & 386 & 497 \\
Score-Acc & 100\% & 100\% \\
\midrule
Key reasoning &
  Coarse step blocks: discriminant proof then Vieta system for part (2). &
  Same rubric, but splits work into smaller blocks (extra granularity in JSON). \\
\bottomrule
\end{tabular}%
}
\end{table}
